%% kapa: exam=1 round=32 year=2021 period=1 form=A subject=real_estate subject_ko=부동산학원론 seq=20 number=100
%% source: Q-Net 감정평가사 기출문제 (cst00309) artlSeq=5208397 file=extracted/1008/5208397/A형/2021년도 제32회 감정평가사 1차 1교시 A형.pdf
%% conversion: kordoc 4.24.51 -> tools/gen1.py (auto); review=proofread-sonnet; flags=none
\begin{question}{100}
부동산 금융에 관한 설명으로 옳은 것은?
\begin{choices}
  \item 역모기지(reverse mortgage)는 시간이 지남에 따라 대출잔액이 늘어나는 구조이고, 일반적으로 비소구형 대출이다.
  \item 가치상승공유형대출(SAM: Shared Appreciation Mortgage)은 담보물의 가치상승 일부분을 대출자가 사전약정에 의해 차입자에게 이전하기로 하는 조건의 대출이다.
  \item 기업의 구조조정을 촉진하기 위하여 기업구조조정 부동산투자회사에 대하여는 현물 출자, 자산구성, 최저자본금을 제한하는 규정이 없다.
  \item 부채금융은 대출이나 회사채 발행 등을 통해 타인자본을 조달하는 방법으로서 저당담보부증권(MBS), 조인트벤처(joint venture) 등이 있다.
  \item 우리나라의 공적보증형태 역모기지제도로 현재 주택연금, 농지연금, 산지연금이 시행되고 있다.
\end{choices}
\end{question}
\answer{1}
